%% kapa: exam=1 round=24 year=2013 period=2 form=A subject=accounting subject_ko=회계학 seq=34 number=34
%% source: Q-Net 감정평가사 기출문제 (cst00309) artlSeq=5073802 file=extracted/1008/5073802/2013년 제24회 감정평가사 1차시험 기출문제/2교시(A형).hwp
%% conversion: kordoc 4.24.51 -> tools/gen1.py (auto); review=proofread-sonnet; flags=none
\begin{question}{34}
(주)감평은 제조원가 항목을 직접재료원가, 직접노무원가 및 제조간접원가로 분류한 후, 개별-정상원가계산을 적용하고 있다. 기초재공품(작업 No.23)의 원가는 ￦22,500이며, 당기에 개별 작업별로 발생된 직접재료원가와 직접노무원가를 다음과 같이 집계하였다. \par
{\small\centering
\begin{tabular}{|c|c|c|}
\hline
\makecell{작업번호} & \makecell{직접재료원가} & \makecell{직접노무원가} \\
\hline
\makecell{No.23} & \makecell{￦2,000} & \makecell{￦6,000} \\
\hline
\makecell{No.24} & \makecell{9,000} & \makecell{10,000} \\
\hline
\makecell{No.25} & \makecell{14,000} & \makecell{8,000} \\
\hline
\end{tabular}\par}

제조간접원가는 직접노무원가에 비례하여 예정배부한다. 기초에 직접노무원가는 ￦20,000으로 예측되었으며, 제조간접원가는 ￦30,000으로 예측되었다. 기말 현재 진행 중인 작업은 No.25뿐이라고 할 때, 당기제품제조원가는? \par
\begin{choices}
  \item ￦34,000
  \item ￦39,500
  \item ￦56,500
  \item ￦62,000
  \item ￦73,500
\end{choices}
\end{question}
\answer{5}
